% Generated draft: requires mathematical and source review.
\documentclass[11pt]{article}
\usepackage[margin=25mm]{geometry}
\usepackage{fontspec}
\setmainfont{Latin Modern Roman}
\usepackage{amsmath,amssymb,mathtools,unicode-math}
\setmathfont{Latin Modern Math}
\usepackage{enumitem,needspace,xurl,hyperref}
\setlength{\parindent}{0pt}
\setlength{\parskip}{5pt}
\setlength{\emergencystretch}{3em}
\newcommand{\N}{\mathbb{N}}
\newcommand{\Z}{\mathbb{Z}}
\newcommand{\Q}{\mathbb{Q}}
\newcommand{\R}{\mathbb{R}}
\newcommand{\C}{\mathbb{C}}
\newcommand{\F}{\mathbb{F}}
\newcommand{\A}{\mathbb{A}}
\newcommand{\PP}{\mathbb{P}}
\newcommand{\E}{\mathbb{E}}
\newcommand{\eps}{\varepsilon}
\newcommand{\dd}{\,\mathrm{d}}
\begin{document}
\section*{UnsolvedMath ID: 30002658}
\textbf{Source problem code:} catalogue entry
\subsection{Definitions and mathematical statement}
A smooth polarized complex projective variety is a pair \((X,L)\), where \(X\) is a smooth complex projective variety and \(L\) is an ample holomorphic line bundle. Regard \(c_1(L)\) as a real cohomology class in \(H^2(X,\R)\). A K\"ahler form \(\omega\) represents this class if \([\omega]=c_1(L)\); it defines a constant-scalar-curvature K\"ahler metric if its scalar curvature \(S_\omega\) is constant on \(X\).

In the test-configuration formulation, K-stability is expressed through positivity of the Donaldson--Futaki invariant for every admissible nontrivial test configuration of \((X,L)\). Such a configuration is an equivariant flat projective degeneration over \(\A^1_{\C}\), equipped with a relatively ample polarization whose nonzero fibers are isomorphic to \((X,L^{\otimes r})\) for some positive integer \(r\). The supplied record does not specify the admissible class of configurations, the treatment of trivial or product configurations, or the precise convention for K-stability; these choices require confirmation from the source.

\textbf{Supplied Yau--Tian--Donaldson conjecture.} For every smooth polarized complex projective variety \((X,L)\) satisfying the source's finite-automorphism-group hypothesis, is the following equivalence true?
\[
\left(\exists\,\omega\text{ K\"ahler on }X:\ [\omega]=c_1(L)\ \text{and}\ S_\omega\text{ is constant}\right)
\quad\Longleftrightarrow\quad
(X,L)\text{ is K-stable}.
\]
The finite-automorphism-group hypothesis is retained as supplied: the record does not explicitly distinguish the automorphism group of \(X\) from the group of automorphisms preserving the polarization. Consequently, the displayed equivalence records the supplied target but still requires these conventions to become a fully specified assertion. The supplied dated literature assessment also questions whether the extraction identifies an unresolved standalone claim, despite the separately supplied corrected statement and verification label.
\subsection{Short English statement}
The supplied conjecture asks whether, for a smooth polarized complex projective variety with finite automorphism group, the polarization class contains a K\"ahler metric of constant scalar curvature exactly when the polarized variety is K-stable. The intended automorphism group and K-stability conventions require clarification. Its supplied open status is dated curation, not an independently verified mathematical conclusion.
\subsection{Sources}
\begin{itemize}
\item Local registry: supplied UnsolvedMath record, permanent ID \(30002658\), exact source problem code \texttt{catalogue entry}, in the Oberwolfach Reports Open Problems collection. Title: \emph{Yau--Tian--Donaldson Conjecture for Polarized Varieties}. The record attributes the conjecture by this name and supplies no individual proposer.
\item Supplied report citation: \emph{Komplexe Analysis (2015)}, Oberwolfach Reports, pp.~2177--2234, DOI \texttt{10.4171/owr/2014/39}. Supplied source URL and literature-source link, labeled ``original Oberwolfach report'': \url{https://doi.org/10.4171/owr/2014/39}. The record separately gives proposed year \(2014\) and an \texttt{owr-year:2014} tag; the relationship to the citation's \(2015\) date is unresolved.
\item Supplied literature check date: \texttt{2026-08-21}. This is supplied dated curation, not independently verified. Its status is ``open,'' with the qualification that the extraction reportedly consists only of the word ``conjecture'' followed by established theorems and a corollary, and thus does not identify an unresolved standalone claim; the queue rule normalized it to open. The supplied background requests independent verification of time-sensitive status and cited claims.
\item The record separately labels the statement \texttt{corrected\_verified} and reports that it was checked against the official Oberwolfach report to restore hypotheses and notation. This verification is a supplied assertion; no source visit or independent verification is claimed here.
\end{itemize}
\end{document}
